\documentclass[12pt,a4paper]{ctexart}

% === 字体设置 ===
\setCJKmainfont{SimSun}
\setmainfont{Times New Roman}

% === 宏包 ===
\usepackage{amsmath}
\usepackage{amssymb}
\usepackage[version=4]{mhchem}
\usepackage{graphicx}
\usepackage{paracol}
\usepackage{xcolor}
\usepackage{fancyhdr}
\usepackage{geometry}

% === 页面设置 ===
\geometry{
  a4paper,
  left=1.5cm,
  right=1.5cm,
  top=2cm,
  bottom=2cm,
  columnsep=1cm
}

% === 内部辅助：计算 paracol 列宽 ===
\newlength{\colwidthinner}
\newcommand{\calccolwidth}{%
  \setlength{\colwidthinner}{\dimexpr\linewidth-2\fboxsep-2\fboxrule\relax}%
}

% 右栏：修改建议框
\newcommand{\correctionbox}[1]{%
  \calccolwidth
  \colorbox{blue!8}{\parbox{\colwidthinner}{#1}}%
}

% === Pandoc 兼容：非标准 LaTeX 命令 ===
\newcommand{\lt}{{<}}
\newcommand{\gt}{{>}}

% === 页眉页脚 ===
\pagestyle{fancy}
\fancyhf{}
\fancyhead[L]{校对报告}
\fancyhead[R]{\leftmark}
\fancyfoot[C]{\thepage}

\begin{document}

\title{test}
\author{}
\date{}
\maketitle

\begin{paracol}{2}

**一本班1**\\
真空中固定有两个点电荷，负电荷$Q_{1}$位于坐标原点处，正电荷$Q_{2}$位于$x$轴上，$Q_{2}$的电荷量大小为$Q_{1}$的$8$倍。若这两点电荷在$x$轴正半轴的$x=x_{0}$处产生的合电场强度为$0$，则$Q_{1}$、$Q_{2}$相距​（ ）\\
A.       $\sqrt{2}x_{0}$      B.       $(2\sqrt{2}-1)x_{0}$      C.       $2\sqrt{2}x_{0}$      D.       $(2\sqrt{2}+1)x_{0}$\\
答案:\\
B\\
解答：                             设$Q_{1}=-Q$，$Q_{2}=8Q$，$Q_{1}$与$Q_{2}$相距$L$，\\
两点电荷在$x$轴正半轴的$x=x_{0}$处产生的合电场强度为$0$，由于$Q_{1} \lt Q_{2}$，故该处距离$Q_{1}$近，距离$Q_{2}$远，故$Q_{2}$应放在$x$轴负半轴\\
该处合场强为零，则两点电荷在该处场强大小相等，方向相反，故有\\
$\frac{k\cdot 8Q}{(L+{x}_{0}{)}^{2}}=\frac{k\cdot Q}{{x}_{0}^{2}}$\\
$L=(2\sqrt{2}-1)x_0$，故$\mathrm{B}$正确，$\mathrm{ACD}$错误。\\
故选：$\mathrm{B}$。\\

\switchcolumn

\switchcolumn*

\end{paracol}

\end{document}
